%% 邮箱回填夹具 —— 主文件（作者**不带**邮箱，模拟 CAS/elsarticle 风格）
\documentclass[11pt]{article}

\usepackage{amsmath}

\title{Explainability-Guided Partial Privileged Learning: A Class-Aware Weighted Similarity Approach}

\author[aff1]{Bo Liu}
\author[aff1]{Yunlong Sun}
\author[aff2]{Yanshan Xiao}

\affil[aff1]{School of Automation, Guangdong University of Technology, Guangzhou 510006, China}
\affil[aff2]{School of Computer Science, Guangdong University of Technology, Guangzhou 510006, China}

\keywords{partial privileged learning, explainable AI, SVM+}

\begin{document}
\maketitle

\begin{abstract}
Learning Using Privileged Information (LUPI) enhances classifier performance by
leveraging additional knowledge during training.
\end{abstract}

\section{Introduction}\label{introduction}

Supervised learning has become a dominant paradigm in machine learning, largely
due to the accessibility of labeled data in various domains. In practical
scenarios, beyond standard labels, we often have access to a priori information
associated with the training samples, which provides richer context for model
interpretation \cite{vapnik2009new}. When a subset of samples lacks this
privileged component, standard LUPI methods cannot be directly applied. We
therefore study partial privileged learning and propose a class-aware weighted
similarity strategy that imputes the missing privileged information while
keeping the downstream classifier calibrated. The remainder of this paper is
organised as follows. Section~\ref{method} describes the proposed framework,
Section~\ref{experiments} reports the empirical study, and Section~\ref{conclusion}
concludes. Our contributions are threefold: first, a formulation of partial
privileged learning that does not assume complete privileged information;
second, an explainability-guided mask estimator that separates discriminative
features from background noise; third, a weighted similarity transfer that
assigns confidence to each imputed privileged vector before it reaches the
classifier. Together these components form a pipeline that is robust to
severe privileged-information scarcity while remaining compatible with
existing SVM+ solvers and requiring no additional annotation effort from
domain experts in the loop.

\section{Methodology}\label{method}

We denote the training set by $\{(x_i, y_i)\}$ and the privileged features by
$x_i^{*}$. The weighted cosine similarity between a query sample and a
candidate sample is computed with respect to the class-conditional importance
vector. Features with high discriminative power retain their weights, while
irrelevant or noisy features are strictly suppressed. The resulting confidence
coefficient modulates the influence of each imputed privileged vector on the
SVM+ objective.

\section{Experiments}\label{experiments}

We evaluate the proposed method on Handwritten Digits, Caltech-101 and
Animals with Attributes 2. Compared with representative baselines, the method
improves classification accuracy and reduces the degradation caused by
privileged-information scarcity.

\section{Conclusion}\label{conclusion}

We presented an explainability-guided framework for partial privileged
learning. Future work will extend the framework to multi-modal settings.

\begin{thebibliography}{9}
\bibitem{vapnik2009new} V. Vapnik and A. Vashist. A new learning paradigm: learning using privileged information. Neural Networks, 2009.
\end{thebibliography}

\end{document}
